\chapter{Operating the product without losing its evidence}
\label{ch:operation}

\section{The bank remains responsible for the control}

The workbench can organize evidence, expose alternatives, verify consequences and
prevent certain unauthorized transitions. It cannot assume the bank's legal or
managerial responsibility. The operational design must identify who owns source
applicability, interpretation, factual mappings, software correctness and release.
A system that assigns every difficult question to a generic ``human review'' queue
has not completed that design.

Meaning review belongs to people authorized to assess the relevant regulatory
requirement. Data owners establish what evidence supports facts, how long it
remains valid and how corrections are represented. Engineering owners establish
that the software preserves the approved specification. Release owners verify the
package, approvals and operational readiness. Business owners determine the
workflow response to unknown or blocked results within the bank's approved policy.

These responsibilities can overlap institutionally, but conflicts and separation
requirements must be explicit. An author should not be able to approve its own
interpretation through a second display name. A model-generated reviewer identity
has no authority. The local prototype's test roles exercise the mechanics; the
bank deployment must bind them to enterprise identity and actual delegated
responsibilities.

The product's value is not that it eliminates judgment. It makes the judgment
specific and attributable. Instead of asking a reviewer to bless a long generated
summary, it asks which definition applies, why an exception attaches to a given
condition, or whether a source is current for the activity. The final control
then carries those decisions into its executable form and future amendment review.

\section{Start with a narrow operational perimeter}

A first deployment should name the activity, legal entity, products, client or
counterparty categories, source families and decision points it supports. This
perimeter is not a marketing description. It determines which circulars enter the
workbench, which factual schemas are required and which host workflows can consume
the resulting package.

For example, a paragraph-10 incentive-control pilot could focus on reviewing a
single component of a proposed offer before it is launched. It would need a
reliable definition of the component, evidence of its product relationship and a
reviewed classification process for gift and fee discount. It would not by itself
cover every marketing, disclosure or distribution requirement. The host must
continue to route the offer through other applicable controls.

Narrowing the perimeter is useful when it removes an unsupported ambiguity by an
explicit operational restriction. If the bank excludes mixed offers from the
automated pilot and routes them to manual review, the system can assess simpler
cases under a clearer subject definition. That restriction must be enforced by
input validation and workflow routing, not merely stated in a project memo.
Otherwise mixed offers will eventually arrive and receive an unintended answer.

The perimeter also defines negative tests. Inputs from another actor, product type
or jurisdiction should produce the specified out-of-scope or unsupported-profile
result. They should not be coerced into the closest known category. A controlled
failure at the boundary is preferable to a plausible answer about the wrong
activity.

Expansion requires evidence. Adding a new circular family may introduce temporal
obligations, different authority relationships or vocabulary that the original
schema does not support. The team should create a new source and method review,
update the factual mappings, and evaluate the expanded task. Success on a narrow
pilot is a reason to consider expansion, not proof that transfer is automatic.

\section{Keep regulatory meaning distinct from bank policy}

A bank can choose a more restrictive operational policy than the minimum condition
identified in a circular. It may decline an uncertain promotion, require an
additional approval or retain a manual check. These decisions should be represented
as bank policy with their own owner, rationale and effective interval. They should
not be folded into the regulatory rule and later cited as if the SFC required them.

This separation improves both accountability and maintenance. If a regulatory
interpretation changes, the team can see whether the bank policy still makes sense.
If a business owner changes an internal approval threshold, it does not appear
that the regulator amended the circular. Reports can explain a blocked offer as
blocked by bank policy pending classification, rather than falsely declaring it
prohibited by the source.

The runtime can evaluate both kinds of controls, provided the bundles identify
their authority and purpose. The host's orchestration specifies how their results
combine. A regulatory condition returning false may still be blocked by a bank
policy; a bank policy returning permissive cannot override a regulatory
prohibition. The combination rules need their own reviewed semantics and tests.

An unresolved interpretation should not be hidden in a conservative default.
Suppose the team treats every cashback as a gift until clarification. That can be
a cautious temporary policy, but the uncertainty report should still state that
the legal classification remains unresolved. Otherwise future reviewers may
mistake a risk-management choice for settled authority and propagate it into
other products.

Temporary policies need expiry or review triggers. An indefinite ``pending
clarification'' rule can become permanent through neglect. The workbench should
record the question, owner, next review condition and affected activities. A
reminder is an operational aid; it does not automatically extend the policy or
resolve the underlying source issue.

\section{Shadow operation is a study, not a hidden release}

Shadow operation runs the proposed control alongside the existing process without
letting its result directly determine the transaction. It allows the team to
observe real input quality, classification workload, latency and disagreement.
The scope and access to real data require the bank's normal authorization.
Shadow results must be clearly distinguished from approved production decisions.

A useful shadow study records the existing process's decision, the workbench's
result, the facts available to each and the timing of those facts. When they
disagree, an independent review examines the reasons. The existing process is not
automatically the truth oracle; manual procedures can be wrong. Equally, the
workbench is not correct merely because it produces a more detailed explanation.
The study needs an adjudication process consistent with Chapter~\ref{ch:evaluation}.

The team should predeclare what would justify promotion, what would veto it and
what would trigger a repair. A severe false clean result might block promotion
while motivating a targeted fix. A data-mapping defect might invalidate a subset
of observations until corrected. A provider outage may reveal operational
fragility without answering the interpretation question. These distinctions keep
the study from becoming a succession of improvised demonstrations.

Shadow operation also reveals whether the selected unit of assessment matches
business practice. If staff enter a whole campaign while the rule expects one
benefit component, the system may receive apparently valid but semantically wrong
inputs. The remedy may be a workflow redesign, not a more capable model. The
interface should help staff identify and evidence components consistently.

Reports from shadow operation remain subject to access, retention and correction
rules. They can contain sensitive facts and preliminary judgments. A later
adjudication should append a correction rather than rewrite the original
observation. This preserves the study's evidence and prevents an apparently
perfect retrospective record from being created by editing away disagreements.

\section{Source monitoring begins a new interpretation cycle}

The workbench should monitor the source families in its approved perimeter for
new publications, revisions and incorporated-material changes. Discovery creates
a candidate source record; it does not automatically establish current
applicability or overwrite the active rule. The acquisition process retains bytes,
metadata, retrieval time and the relationship to prior versions.

A textual difference is only the first signal. A changed definition may affect
many controls even when their operative paragraphs are unchanged. A formatting
change may have no substantive consequence but still change the file hash. The
impact service therefore combines byte identity, extracted text differences,
source-unit mapping and the explicit dependency graph.

The first impact report identifies which candidates, facts, rules, duties,
reports and releases depend on the changed material. It should distinguish direct
citations from indirect definition dependencies. A control can be affected through
an incorporated code provision that is not quoted in its main explanation. This
is why dependency closure was required before initial interpretation.

An unchanged structural dependency report is not a proof of semantic equivalence.
Conversely, a changed paragraph number does not prove a changed obligation.
Semantic comparison and legal review are separate actions. The ensemble can
propose interpretations of the change and construct scenarios on which old and
new rules differ, but a reviewer must establish the intended effective regime.

Every changed rule version requires the relevant review, build and verification
steps again. Reusing a previously approved JAR because its human-readable title
looks the same would break the evidence chain. Reusing a verified runtime library
can be legitimate if its identity and supported semantics remain unchanged, while
the new policy bundle and generated binding receive fresh evidence.

\section{A worked amendment scenario}

Consider a synthetic amendment used only to explain the process. The original
selected rule contains a fee-discount exception. A later source adds a condition
requiring a particular relationship between the discount and the transaction.
This is not a claim that the SFC has made such an amendment; it is a controlled
example of how a change would propagate.

The source custodian retains the new document and identifies the added condition.
The inventory maps it to the earlier exception and opens an interpretation run.
Candidate A treats the condition as applying to all offers assessed after the
new commencement date. Candidate B distinguishes offers already committed before
that date under a possible transition provision. The source packet must include
that transition provision if one exists; the system cannot infer grandfathering
from business preference.

The scenario designer creates offers before and after commencement, an offer
agreed before but performed after, and an offer with incomplete timing evidence.
These examples reveal whether the competing readings differ. The formal checker
can compare their executable conditions once the timing concepts are defined.
It cannot determine which transition interpretation the source supports without
the relevant authority.

When reviewers accept the new interpretation, the data owner confirms that the
required relationship and timing facts can be supplied. Engineering builds and
verifies the new Java package. The release manifest binds its effective interval
and applicability profile. The host's release selector is tested at the exact
boundary, including the excluded endpoint of the earlier interval.

Existing event obligations require separate treatment. If a duty was opened under
the old bundle, the current MVP preserves that binding. A legal decision may
require migration, continuation or another treatment. The workbench records the
question and any approved migration specification. No background update silently
changes the source version attached to a historical duty.

The old source, interpretation, JAR and decisions remain available for replay.
A later investigation can ask what the system would have decided under the old
rule with the evidence known then, or what a corrected assessment under the new
rule would conclude. Those are different questions and their results carry
different assessment contexts.

\section{Bitemporal replay explains rather than rewrites history}

The two-time model is particularly useful in an incident. Valid time asks when a
fact or rule applies to the underlying event. Knowledge time asks what information
was available to the system at a specified point. A replay must state both.
Without this distinction, later corrections can make earlier decisions appear
inexplicably wrong or retrospectively perfect.

Suppose a transaction was assessed on Monday using a classification record
available that morning. On Wednesday, a reviewer finds evidence that changes the
classification effective from Monday. A historical replay with Monday knowledge
reproduces the original evidence state. A corrected assessment with Wednesday
knowledge can produce a different result for the same Monday transaction. Both
results are legitimate records of different questions.

The system should label the corrected assessment and link it to the original.
It should not replace the original execution hash or pretend that Wednesday's
evidence was available on Monday. This distinction supports fair investigation
of whether the failure concerned the rule, data, process or delayed information.
It also identifies which other decisions relied on the same evidence.

Replay requires the exact historical runtime or a verified equivalent path, the
bundle, fact snapshot, event history where relevant, calendars and assessment
context. Retaining only a final status and a current source URL is insufficient.
The retention policy must preserve these dependencies for the bank's required
period; this book does not invent a statutory retention duration.

Reproducing an old decision does not endorse it. A replay can faithfully show that
an earlier rule omitted an exception. That is useful evidence about the incident.
The corrected rule and reassessment should be linked, with the reason for change
and any operational remediation. Traceability supports accountability precisely
because it can preserve an error honestly.

\section{Operational responses to unknown and conflict}

Unknown and conflict are ordinary outcomes in a system built around evidence.
They should have designed workflows rather than being treated as exceptional
software crashes. The host policy identifies the responsible queue, information
request, permitted temporary action and escalation deadline for each reason.
These choices depend on the bank's process and must be approved there.

A missing classification might route to a product specialist. Conflicting fee
records might route to the data owner. An unavailable incorporated source might
require legal review before any automated decision. An unsupported RuleIR operator
requires engineering and meaning review together. The reason code helps select
the remedy, but the workflow should expose the substantive question in ordinary
language.

A conflict should preserve both records and their provenance. Automatically taking
the newest record can hide a disagreement between systems or an unauthorized
correction. If the bank adopts a precedence rule, that rule must be explicit and
reviewed for the evidence type. It should not be embedded in a generic database
query that happens to order by timestamp.

Manual overrides need their own authority and scope. An override may authorize a
specific operational action under bank policy while leaving the regulatory
interpretation unresolved. The record should say exactly what was overridden,
by whom, for which case and period, and with what rationale. It must not change
the underlying source or make future cases automatically inherit the exception.

Queues themselves can fail. Unreviewed items can accumulate, deadlines can pass,
and staff can misunderstand a blocked status. Monitoring should therefore track
age, owner, materiality and next action for unresolved issues. A low model error
rate is little comfort if the operational process routinely ignores the issues
that the model correctly escalates.

\section{Incident analysis should locate the failed boundary}

When an incorrect control or decision is discovered, the first question is where
the evidence chain failed. Was the wrong source acquired? Was a qualification
missed? Was the accepted interpretation encoded incorrectly? Did the host supply
misclassified facts? Did a transaction race invalidate the snapshot? Did an
unauthorized or stale package reach production? These mechanisms require different
repairs and different assessments of affected decisions.

The incident record should preserve the exact failing example, source versions,
candidate, approvals, JAR, input snapshot, result and operational action. It
should also preserve the strongest competing interpretation that was available
at the time. If the objection was present but suppressed, the remedy concerns
reporting and review as well as generation quality.

Containment is a bank decision. It may involve suspending a particular offer
category, returning to an approved manual process, disabling a faulty adapter or
replacing an implementation while retaining the same accepted rule. The incident
team must distinguish these options from reverting to a legally superseded rule.
The workbench supplies impact information; it does not decide the bank's external
communications or regulatory response.

Root-cause analysis should avoid a generic conclusion that ``the AI hallucinated.''
That phrase can conceal a missing source, a contaminated oracle, an undefined
classification or a budget policy that forced premature closure. A specific
mechanism supports a specific preventive control. The discovered failure becomes
a retained regression scenario, with its source and interpretation status clear.

The team should also ask what similar decisions may be affected. Dependency and
lineage queries can identify every package using the disputed definition and
every assessment using the relevant fact mapping. A database search by final
status alone will miss decisions that happened to be correct for other reasons.
The scope of remediation follows the causal dependency, not merely the visible
outcome.

\section{Availability failures must not change the meaning of a result}

The production Java evaluator should not depend on a live model or source website.
That separation lets an approved control continue to execute during an
interpretation-provider outage, subject to the bank's release and source-monitoring
policy. An outage in the workbench still matters for new controls and amendment
review, but it need not make every existing decision nondeterministic.

The host must distinguish an unavailable interpretation service from an unknown
fact and from an expired release. These conditions can require different actions.
An approved, still-applicable local package may remain usable when a research
provider is unavailable. A release whose effective status cannot be established
presents a different problem. The fallback policy must state the conditions under
which each path is permitted.

No fallback may convert a technical error into a substantive false result. If a
Java evaluation fails because the snapshot is malformed, returning false would
imply that the selected condition was evaluated and not satisfied. The correct
result is an error with an operational disposition. Likewise, a missing
interpretation report cannot be replaced by a default approval flag.

The ensemble controller has its own recovery obligations. A restart recovers
reserved and dispatched actions, preserves their charges and determines which
results are known, late or unknown. It does not restart the investigation from
zero unless an authorized successor run is created. The supervisor monitors
absolute deadlines independently of the worker's progress messages.

Capacity planning should consider the difficult tail: large dependency graphs,
many genuine alternatives and long-running external checks. Graph and action
limits protect the service, but they create explicit incomplete outcomes. The
operational owner needs to know how often those outcomes occur and whether the
manual review process can absorb them. Increasing limits without measuring cost
and review quality is not a complete capacity strategy.

\section{Model updates are method changes}

A new model version can change which alternatives are generated, how sources are
summarized and which questions are prioritized. Even if the JSON schema remains
unchanged, the interpretation method has changed. The workbench should bind
provider, model, prompt and configuration versions to each run and evaluate
material updates before making them the default.

The update review asks whether known omission cases remain detected, minority
readings remain visible, unsupported alternatives increase, and uncertainty scores
retain their measured relationship to the reference. It also checks provider
behavior relevant to budgets, such as retries, tool use and token accounting.
A faster model that silently changes these behaviors can invalidate the
controller's operational assumptions.

Existing approved Java packages do not need to be regenerated merely because the
proposal model changes. Their authority derives from the retained interpretation,
review and verification evidence. A new model may identify a concern that triggers
review, but its existence alone does not invalidate every prior decision. This
separation limits unnecessary churn while allowing genuine discoveries to enter
an accountable amendment or incident process.

Prompt changes deserve similar treatment. Adding a leading example can improve
performance on one family while anchoring all members to the same interpretation.
Removing a request for counterarguments can reduce cost while weakening diversity.
The change record should explain the intended mechanism, evaluation evidence and
remaining uncertainty. A prompt is part of the method, not incidental prose.

Search-policy changes also require evidence. Switching from a deterministic
scheduler to MCTS, changing family quotas or increasing a no-progress threshold
alters which questions receive attention. The study should compare these choices
under equal resources and independent judgments. Passing controller invariants
is necessary for either scheduler, but does not establish which investigates law
more effectively.

\section{A release readiness review with concrete evidence}

A readiness review should inspect a concrete package, not approve a future promise
to generate one. The package contains the accepted interpretation dossier, source
coverage, factual mapping, verification report, build and release manifests,
operational dispositions and unresolved-risk record. All identities must agree.
An approval of one candidate cannot be reused for a different generated bundle.

The meaning reviewer checks that the selected source perimeter and assumptions
support the rule, that plausible alternatives have a reasoned disposition, and
that any remaining uncertainty is compatible with the proposed scope. The data
owner checks input definitions, evidence quality, freshness, corrections and
calendar or valuation policies. Engineering checks exact behavior, traces,
resource limits and integration. The release owner checks authority, deployment
mechanics, monitoring and recovery.

A useful readiness table separates evidence from a decision:
\begin{longtable}{@{}P{31mm}P{60mm}P{57mm}@{}}
\caption{Evidence required for a bounded deployment decision.}\\
\toprule
Area & Evidence to inspect & Question that remains a human responsibility\\\midrule\endfirsthead
\toprule Area & Evidence to inspect & Question that remains a human responsibility\\\midrule\endhead
Meaning & Source inventory, candidates, objections, adjudication & Is this reading justified for the selected activity and time?\\
Data & Fact dictionary, provenance, validity, conflict policy & Do the bank's records establish the facts the rule needs?\\
Execution & Conformance, witnesses, mutations, Java comparison & Is the supported implementation scope adequate for the process?\\
Authority & Exact manifests, reviewer roles, signatures & Are the required people authorized and accountable for this release?\\
Operation & Host dispositions, race tests, fallback and incident rehearsal & Can staff act correctly when the result is unknown or blocked?\\
Evidence & Retained sources, histories, reproducible reports & Can a later reviewer reconstruct the decision without relying on memory?\\
\bottomrule
\end{longtable}

A missing item should produce a specific block and next action. It should not
trigger a vague demand for more confidence. If a classification is unresolved,
obtain the authority or narrow the scope. If Java traces differ, repair the
implementation. If enterprise identity is not integrated, complete that integration
before relying on local synthetic roles. The package makes these distinctions
visible enough to assign work.

A review may approve a limited deployment with explicit exclusions and monitoring
conditions. Those conditions become enforceable configuration and expiry or
review triggers. They must not remain only in meeting minutes. The system should
be able to show which operational path enforces each exclusion and which test
would fail if it were bypassed.

\section{Three operational cases for the implementation team}

The first case is a missing definition. The source packet is intact, the candidate
formula compiles, all named scenarios pass, and the ensemble agrees on the
expression. An incorporated definition is unavailable. The correct report is
blocked for the affected classification, even though the engineering checks pass.
The next action is to obtain or adjudicate the definition, not to add more random
truth-table tests.

The second case is a wrong host mapping. Reviewers accepted a component-level
rule, but the host sends one Boolean classification for an entire mixed offer.
The Java evaluator faithfully applies its inputs and returns a reproducible
result. The error lies at the subject and fact boundary. The remedy includes a
component-aware adapter, affected-case analysis and a regression scenario that
would expose a whole-offer mapping. Reproving the Boolean evaluator alone would
not address the failure.

The third case is an unresolved exception priority. Two exception guards are
true. The current RuleIR default returns conflict, while a developer proposes to
choose the first exception in the array. That change would make the program run,
but it would introduce a priority that the specification did not authorize. The
workbench must preserve conflict and ask whether the source establishes priority,
whether explicit nesting is justified, or whether the rule lies outside the
supported fragment.

These cases demonstrate why completion criteria must follow the claimed target.
In the first, source evidence is incomplete. In the second, the application mapping
is wrong relative to the accepted interpretation. In the third, a proposed
implementation change would alter semantics. The same dashboard label ``needs
review'' is not enough; the system must name the responsible boundary and the
evidence needed to repair it.

A team rehearsal should walk through each case using the actual interface and
reports. Staff should be able to locate the source, identify the affected
candidate, understand the block and record a decision without editing hidden
configuration. Rehearsal findings become product changes and acceptance tests.
A successful technical build does not substitute for this operational exercise.

\section{Maintaining a useful literature and method record}

The paper library is a working research asset. Each retained paper has a stable
filename, source URL, version or publication information, digest and reading
status. The title-and-first-author naming convention requested for this project
helps people locate a paper, while the manifest disambiguates revised preprints
and related editions. A citation should point to the version actually inspected.

The bibliography and reading notes distinguish published results, local
derivations and proposed adaptations. Catala's treatment of defaults, Stipula's
state and resource semantics, argumentation frameworks, tree search and uncertainty
methods answer different questions. Combining them in one architecture does not
transfer every theorem or empirical result to the combination. The method record
states which mechanism is borrowed and what must be re-established locally.

Official code is valuable evidence about practical assumptions. The semantic-
clustering example in Chapter~\ref{ch:search} shows why code and paper can require
separate inspection. A pinned repository revision may implement a broader
condition than the prose suggests. The correct response is to document and test
the difference, not silently choose the version that makes the proposal look
stronger.

The literature record should also preserve inaccessible or uninspected leads.
They can guide future work, but cannot support categorical implementation claims.
The present review is a substantial bounded survey of the retained sources,
not a declaration that no relevant paper exists outside the library. New work can
change the best engineering choice, especially in a rapidly developing area.

ResearchAssistant is useful for acquisition, parsing and source structure.
MathDevMCP is useful for named formal obligations. DynareMCP contributes patterns
for explicit hypotheses, falsifiers and bounded investigations. Their reuse
should remain concrete and versioned. The bank should be able to identify which
function, data structure or proof obligation is being borrowed, rather than
relying on a broad claim that another project already solved the difficult part.

\section{A practical sequence from manuscript to product}

The first step is contract implementation with public sources and scripted
members. This establishes that the ensemble can preserve alternatives, detect
specified discrepancies, respect limits and produce complete reports. The
second-circular example provides a running integration case, while synthetic
fully resolved cases exercise the positive review route. The original blocked
example should remain blocked unless actual new evidence resolves its issues.

The second step is independent reference construction and restricted live-model
evaluation. The team freezes source inventories and judgments before seeing
model outputs, runs the baseline ladder and records uncertainty. Failures are
classified by boundary and used for planned repairs. No model is promoted because
it merely reproduces the author's original 22-case oracle.

The third step is enterprise integration: identity, approved data access, isolated
source parsing, provider controls, signed or equivalently authenticated releases,
monitoring and host transaction behavior. This work can proceed alongside method
evaluation where dependencies permit. It should not be postponed until a final
demo, because operational constraints may change which designs are feasible.

The fourth step is authorized shadow operation in a narrow activity perimeter.
The study compares the workbench with the existing process and independently
adjudicates disagreements. It measures both substantive quality and reviewer
workload. The team rehearses source amendments, unknown facts, provider outages,
worker crashes and rollback decisions.

The fifth step is a bounded deployment decision based on the concrete evidence
package. The decision may approve only a subset of controls or retain manual
review for disputed categories. Monitoring then continues the evidence process:
new sources, unexpected disagreements and incidents create successor
investigations. Deployment is not the end of interpretation assurance.

This sequence does not assign an arbitrary calendar duration or staffing number.
The bank's source perimeter, access constraints and review capacity determine
those estimates. The task dependencies and acceptance evidence are concrete enough
to build a schedule once those inputs are known. Inventing a confident delivery
date without them would add apparent precision without resolving the underlying
work.

\section{What the completed prototype will and will not establish}

A completed deterministic prototype will establish that the specified records,
controller, budgets, reports and release checks behave as tested. It will produce
actual Java packages from accepted RuleIR and preserve the existing deterministic
execution contract. It will demonstrate that known omissions and persistent
ambiguities lead to the intended different outcomes in scripted scenarios.

A successful independent interpretation study can establish empirical performance
on its declared task population, with uncertainty and limitations. It may support
a bounded use or a choice among methods. It cannot prove that every future
circular will be interpreted correctly, especially under new vocabulary,
authority or distribution shift. A formal proof of the controller's release
invariant similarly depends on the recorded evidence and does not prove that no
source obligation was ever missed.

These limits do not make the product pointless. The present manual process also
contains interpretation, transcription, coding and review errors, often without a
complete record of competing readings. The proposed workbench makes those stages
explicit, introduces complementary checks and preserves the questions that remain.
Its value should be demonstrated against that real process with a fair reference
and comparable resources.

The high stakes justify both ambition and precision. The system should actively
search for ways its preferred reading could be wrong, retain those challenges,
and stop automatic work honestly when the evidence or budget is insufficient.
It should deliver code whose relationship to an approved specification can be
checked, while keeping the approval of meaning visible as a separate act.
That combination is a concrete engineering objective and a defensible basis for
building the next version of LegalMath.
